\documentclass[11pt]{article}
\usepackage[T1]{fontenc}
\usepackage[utf8]{inputenc}
\usepackage{lmodern}
\usepackage[english]{babel}
\usepackage{amsmath,amssymb,mathtools,array,booktabs,pdflscape,adjustbox,diagbox}
\usepackage{geometry}
\usepackage{microtype}
\geometry{a4paper,margin=1.45cm}
\setlength{\parindent}{1.25em}
\setlength{\parskip}{0pt}
\makeatletter
\renewcommand\footnotesize{\@setfontsize\footnotesize{7.5}{8.5}\selectfont}
\makeatother
\newcommand{\widebar}[1]{\overline{#1}}
\newcommand{\A}{\Delta}
\newcommand{\D}{\Delta}
\allowdisplaybreaks
\setlength{\emergencystretch}{10em}
\sloppy
\hbadness=10000
\vbadness=10000
\hfuzz=2pt
\vfuzz=10pt
\raggedbottom
\newcommand{\R}{\varrho}
\newcommand{\hata}{\widehat{a}}
\newcommand{\hatv}{\widehat{v}}
\newcommand{\modu}[1]{\;\operatorname{mod} #1}
\newcommand{\mcell}[1]{\ensuremath{\begin{gathered}#1\end{gathered}}}
\providecommand{\pair}[2]{(#1\,|\,#2)}
\providecommand{\eps}{\varepsilon}
\providecommand{\rhoidx}{\varrho}

% Batch 20 macros
\providecommand{\dd}{\mathrm d}
\providecommand{\Div}{\operatorname{Div}}
\providecommand{\G}{\mathfrak G}
\providecommand{\bd}{\bar{\delta}}
\providecommand{\p}{\partial}


% Batch 21 ideal/function-field macros
\providecommand{\frS}{\mathfrak S}
\providecommand{\frp}{\mathfrak p}
\providecommand{\fra}{\mathfrak a}
\providecommand{\frb}{\mathfrak b}
\providecommand{\frc}{\mathfrak c}
\providecommand{\frf}{\mathfrak f}
\providecommand{\frr}{\mathfrak r}
\providecommand{\frm}{\mathfrak m}
\providecommand{\frD}{\mathfrak D}
\providecommand{\calA}{\mathcal A}
\providecommand{\calB}{\mathcal B}
\providecommand{\calN}{\mathcal N}


\providecommand{\p}{\partial}
\providecommand{\dd}{\mathrm d}
\providecommand{\modu}[1]{\;\operatorname{mod} #1}
\newcommand{\mM}{\mathfrak{M}}
\newcommand{\mN}{\mathfrak{N}}
\newcommand{\mL}{\mathfrak{L}}
\newcommand{\mG}{\mathfrak{G}}
\newcommand{\mA}{\mathfrak{A}}
\newcommand{\mB}{\mathfrak{B}}
\newcommand{\mX}{\mathfrak{X}}
\newcommand{\mY}{\mathfrak{Y}}
\newcommand{\ma}{\mathfrak{a}}
\newcommand{\mb}{\mathfrak{b}}
\newcommand{\mx}{\mathfrak{x}}
\newcommand{\my}{\mathfrak{y}}
\newcommand{\mz}{\mathfrak{z}}
\newcommand{\me}{\mathfrak{e}}
\newcommand{\mbarA}{\overline{\mathfrak{A}}}
\newcommand{\mbarB}{\overline{\mathfrak{B}}}
\newcommand{\bara}{\overline{\mathfrak{a}}}

\begin{document}
\clearpage
\setcounter{footnote}{0}
\section*{17. Together with W. Schmeidler: Modules in Noncommutative Domains, Especially from Difference Expressions}
\begin{center}
Math. Zs. 8 (1920), pp. 1--35
\end{center}

\begin{center}
\textbf{Contents.}
\end{center}
\begin{center}
\begin{minipage}{0.82\linewidth}
Introduction.\\
\S~1. Formal matters concerning differential and difference expressions.\\
\S~2. The general polynomial domain.\\
\S~3. One-sided modules and their residue groups.\\
\S~4. Reducibility of residue groups into two subgroups and their relation to the module.\\
\S~5. Computation with units and reducibility into $k$ subgroups.\\
\S~6. Finiteness theorems.\\
\S~7. A case of unique decomposition.\\
\S~8. Additive decomposition of completely reducible groups into prime groups. The theorem on isomorphism.\\
\S~9. Connection between the concepts ``isomorphic'' and ``of the same kind''.\\
\S~10. Existence of infinitely many decompositions of a group.\\
\S~11. Relations to systems of differential equations and their integrals.\\
\S~12. Examples.
\end{minipage}
\end{center}

\begin{center}\textbf{Introduction.}\end{center}

The formal theory of differential expressions in one variable leads to a decomposition into irreducible factors which is unique in a certain sense.\footnote{Cf. E. Landau, Ein Satz über die Zerlegung homogener linearer Differentialausdrücke in irreduzible Faktoren, Journal für die reine und angewandte Mathematik 124 (1902), pp. 115--120, and subsequently A. Loewy, Über reduzible lineare homogene Differentialgleichungen, Mathematische Annalen 56 (1903), pp. 549--584.} But whereas the corresponding theorem in algebra can be carried over to polynomials in several variables, this is no longer possible for partial differential expressions.\footnote{Cf. a counterexample of Landau published in H. Blumberg, Über algebraische Eigenschaften von linearen homogenen Differentialausdrücken, dissertation, Göttingen 1912, 52 pp., pp. 51--52.} For differential expressions in one variable, however, alongside the product representation there is also a representation as least common multiple, with similar laws holding; for, in two different decompositions, the individual constituents can be assigned to one another pairwise: they are ``of the same kind''.\footnote{Cf., besides the papers just cited, A. Loewy, Über vollständig reduzible lineare homogene Differentialgleichungen, Mathematische Annalen 62 (1906), pp. 89--117; cited as Loewy.}

This concept of the least common multiple can now be interpreted as a module concept and so transferred to $n$ variables; in the present paper this gives the natural generalization of the theorems just mentioned. Beyond this, however, we also note that for differential expressions only the following fact is used: they can be interpreted as polynomials with noncommutative multiplication, for which the degree of a product is equal to the sum of the degrees of the factors. These conditions are fulfilled, for example, by difference expressions as well, and hence the theorems also hold for these, where until now they were likewise known only for one variable.\footnote{Cf. G. Wallenberg--A. Guldberg, Theorie der linearen Differenzengleichungen, Leipzig and Berlin 1911.}

We therefore develop, in general, a theory of modules whose elements are polynomials with noncommutative multiplication; in doing so we reduce the study of the modules to that of residue classes.\footnote{For the case of commutative domains, cf. W. Schmeidler, Über Moduln und Gruppen hyperkomplexer Größen, Mathematische Zeitschrift 3 (1919), pp. 29--42.} In contrast with the commutative special case, one can indeed speak here of the sum of residue classes, but not of their product; one can speak only of the product of a residue class with a polynomial. The concept of a ``residue group'' is therefore to be modified relative to the commutative case. In the special case of polynomials in one variable, the residue group is moreover ``finite'', that is, it possesses only a finite number of linearly independent residue classes. In general there is no such finite ``order''; the groups are ``infinite''. For our methods, however, this difference is immaterial.

Just as in the commutative case, to the representation of a module as the least common multiple of relatively prime modules there corresponds the simpler process of additive decomposition of the residue group. In this connection the isomorphism of residue groups proves essential; upon specialization to differential expressions in one variable this isomorphism becomes the concept of the same kind for the associated modules, while upon specialization to commutative modules it becomes identity. Incidentally, the concept of the same kind can also be transferred to $n$ variables, and its agreement with the isomorphism concept can be proved. The fact that in the additive decomposition of the residue group only finitely many summands ever occur follows by transferring Hilbert's theorem on the module basis, using Gordan's method of proof; Hilbert's original method of proof would require restrictions on the multiplication laws (in formula (6), $a$ on the right instead of $a_i$, which is fulfilled for differential expressions but not for difference expressions).

In order to arrive at a uniqueness theorem, we must specialize the modules -- as Loewy also does -- to ``completely reducible'' ones, that is, to those representable as least common multiples of prime modules. Given two different decompositions, the constituents are then either pairwise identical, or at least to every constituent of one decomposition there corresponds an isomorphic constituent of the other decomposition, and conversely; at the same time, each decomposition then contains at least two isomorphic constituents. Furthermore, from the occurrence of two isomorphic constituents in one and the same decomposition there follows the existence of infinitely many decompositions; this even holds without the restriction to prime modules, which until now was not known even in the special case of differential expressions in one variable.

Finally, for modules from differential expressions, we discuss the connection with integrals and show that, for finite residue groups, the order is equal to the largest number of linearly independent integrals. Thus arbitrary functions occur in the general integral precisely when the residue group is not finite. The concept that two modules are of the same kind means, here as in the one-variable case, a ``birational'' relation between the integrals. Some examples conclude the paper.

The first impulse for the present paper was given several years ago by a question of E. Landau to E. Noether concerning the generalization of the product theorem. At that time E. Noether only observed that the question was one of module theory in the noncommutative polynomial domain, but could not yet be attacked. The Lasker module theorems then known were not directly transferable to this domain, because Lasker's method of proof rests on the theorem that a polynomial is uniquely decomposable into irreducible factors.\footnote{Recently E. Noether observed that the Lasker theorems too can be established without using this theorem, by methods related to those used here.} Only Schmeidler's methods offered the possibility of a transfer, which we then carried out jointly. In detail, let it be mentioned that the generalization of the residue group, the transfer and application of Hilbert's finiteness theorem, and the existence of infinitely many decompositions are due to E. Noether; the generalization of Loewy's concept of complete reducibility, the concept of isomorphism, and its connection with the concept of the same kind are due to W. Schmeidler.

\subsection*{\S~1. Formal Matters Concerning Differential and Difference Expressions.}

1. For differential expressions in one variable, a symbolic product formation has been introduced. Let
\[
 a_0(x)\frac{\dd^n y}{\dd x^n}+a_1(x)\frac{\dd^{n-1}y}{\dd x^{n-1}}+\cdots+a_n(x)y=A(y),
\]
\[
 b_0(x)\frac{\dd^m y}{\dd x^m}+b_1(x)\frac{\dd^{m-1}y}{\dd x^{m-1}}+\cdots+b_m(x)y=B(y).
\]
Then one defines the symbolic product $AB(y)$ as the differential expression
\[
 a_0(x)\frac{\dd^n B(y)}{\dd x^n}+a_1(x)\frac{\dd^{n-1}B(y)}{\dd x^{n-1}}+\cdots+a_n(x)B(y).
\]
This product formation is in reality a multiplication of two operators; it is most simply written by replacing, as usual, $\frac{\dd^n}{\dd x^n}$ by $\left(\frac{\dd}{\dd x}\right)^n$. Then $AB$ becomes
\[
 \left(a_0\left(\frac{\dd}{\dd x}\right)^n+\cdots+a_n\right)
 \left(b_0\left(\frac{\dd}{\dd x}\right)^m+\cdots+b_m\right)y,
\]
where the ordinary rules for differentiating sums and products apply in carrying out the product. Every such operator can now be regarded as a polynomial in the single variable $\frac{\dd}{\dd x}=\xi$, by not writing explicitly the function to which the operator is applied and by fixing for the variable $\xi$ the multiplication rule corresponding to product differentiation:
\[
 \frac{\dd}{\dd x}\bigl(a(x)y\bigr)=a(x)\frac{\dd}{\dd x}(y)+a'(x)y,
\]
or, in our notation $(n=1,m=0)$,
\begin{equation}
 \xi\cdot a=a\cdot\xi+a'.
\tag{1}
\end{equation}

One proceeds in exactly the same way for linear partial differential expressions
\[
 \left(\sum a_{\alpha_1\ldots\alpha_n}(x_1,\ldots,x_n)
 \frac{\p^{\alpha_1+\cdots+\alpha_n}}
 {\p x_1^{\alpha_1}\cdots\p x_n^{\alpha_n}}\right)y
\]
of a function $y$. For the $n$ operators $\frac{\p}{\p x_1},\ldots,\frac{\p}{\p x_n}$ we introduce variables $\xi_1,
\ldots,\xi_n$ and thereby obtain the polynomial
\[
 \sum a_{\alpha_1\ldots\alpha_n}(x_1,\ldots,x_n)
 \xi_1^{\alpha_1}\cdots\xi_n^{\alpha_n}.
\]
The product of two polynomials is then computed according to the laws
\begin{equation}
 \xi_i a=a\xi_i+a_i \qquad (i=1,\ldots,n),
\tag{2}
\end{equation}
where $a_i$ denotes the partial derivative of $a$ with respect to $x_i$; the products of the $\xi$'s among themselves are commutative. Apart from the commutative law of multiplication, all laws of addition and multiplication hold for the polynomials so defined just as in ordinary algebra. In particular, for the product of two such polynomials the total degree, and the degree in each individual variable, is equal to the sum of the corresponding degrees of the two factors.

2. For difference expressions (cf. Wallenberg, loc. cit.) of the form
\[
 a_x^{(0)}y_{x+n}+a_x^{(1)}y_{x+n-1}+\cdots+a_x^{(n)}y_x=A_x,
\]
\[
 b_x^{(0)}y_{x+m}+b_x^{(1)}y_{x+m-1}+\cdots+b_x^{(m)}y_x=B_x
\]
one has a corresponding product definition
\[
 (AB)_x=a_x^{(0)}B_{x+n}+a_x^{(1)}B_{x+n-1}+\cdots+a_x^{(n)}B_x,
\]
which can be carried out with the aid of the multiplication rule
\[
 (ay)_{x+1}=a_{x+1}y_{x+1}.
\]
If we introduce the operator $\xi$ for the transition from $x$ to $x+1$, then we obtain $\xi\{ay\}=a^*\xi\{y\}$, where $a^*=a_{x+1}$. Again one can omit the function $y$ to which the operator is applied and obtains the law
\begin{equation}
 \xi a=a^*\xi,
\tag{3}
\end{equation}
where, along with $a$, also $a^*$ does not vanish identically. In exactly the same way, for $n$ variables,
\begin{equation}
 \xi_i a=a_i\xi_i\qquad (i=1,\ldots,n),
\tag{4}
\end{equation}
where $\xi_i$ stands for the transition from $x_i$ to $x_i+1$, and $a_i$ denotes the function obtained from $a$ by replacing $x_i$ by $x_i+1$; hence it does not vanish identically if $a$ does not. With these conventions, all rules of addition and multiplication, including the rule concerning the degree of a product, again hold, except for the commutative law.

Both cases will be subordinated, in the following section, to a quite generally defined polynomial domain with arbitrary coefficients.

\subsection*{\S~2. The General Polynomial Domain.}

We start from an arbitrary abstractly defined field $P$ (whose elements we denote by small Latin letters), for which all rules of algebra hold, in particular also the commutative law of multiplication with unique inversion. To this field we adjoin $n$ indeterminates $\xi_1,\ldots,\xi_n$; that is, we consider the integral domain $J$ consisting of the products $a\xi_1^{\nu_1}\xi_2^{\nu_2}\cdots\xi_n^{\nu_n}$ and their finite sums, where all laws of addition and multiplication are to hold with the exception of the commutative law of multiplication. In particular, we call finite sums of the form
\[
 F(\xi_1,\ldots,\xi_n)=\sum a_{\nu_1\ldots\nu_n}\xi_1^{\nu_1}\cdots\xi_n^{\nu_n}
\]
``polynomials'' (which throughout are to be denoted by capital Latin letters). In place of the commutative law, we now subject the elements of $J$ to product laws such that \emph{the product of two polynomials is again a polynomial}; thus the domain of all polynomials from $J$ already exhausts $J$. For this it is plainly necessary and sufficient that, for the elements $\xi_1,
\ldots,\xi_n$ and an arbitrary element $a$ of $P$, product rules of the form
\begin{equation}
 \xi_i a=F_i(\xi_1,\ldots,\xi_n)\qquad (i=1,\ldots,n)
\tag{5}
\end{equation}
hold, where $F_i(\xi_1,\ldots,\xi_n)$ denotes some polynomial depending naturally on $\xi_i$ and $a$. (In particular, the rules (2) and (4) for differential and difference expressions, respectively, are of this form.) For since $\xi_i a$ is not in itself a polynomial, while by hypothesis it belongs to the domain, the equation is necessary. Conversely, by finitely repeated application of formula (5), every element of the domain $J$ can be transformed into a polynomial.

For the unit of the field, which we denote by 1, we shall further require $\xi_i\cdot1=1\cdot\xi_i=\xi_i$. We shall also stipulate that the multiplication of $\xi_1,\ldots,\xi_n$ among themselves is commutative.\footnote{This is needed only from \S~6 onward for the finiteness theorem and its consequences.} Every polynomial in $J$ then writes itself just like an ordinary polynomial in $\xi_1,\ldots,\xi_n$ with coefficients from $P$, and two polynomials agree only when all their coefficients agree.

More specifically, in what follows\footnote{This too is used only from \S~6 onward.} we require that, in place of (5), the rule
\begin{equation}
 \xi_i a=a_i\xi_i+b_i\qquad (a_i\ne0)\quad(i=1,\ldots,n)
\tag{6}
\end{equation}
shall hold. Here too, in particular, the laws (2) and (4) are included. Then the degree of a product is equal to the sum of the degrees of the factors; this holds both for total degree and for the degree in each individual variable, so that the product of two polynomials which do not vanish identically is likewise different from 0.

\subsection*{\S~3. One-Sided Modules and Their Residue Groups.}

For our domain $J$, because multiplication is noncommutative, the following kinds of modules\footnote{We denote modules by capital German letters.} can be distinguished:\footnote{One-sided ideals were already considered by A. Hurwitz; cf. Vorlesungen über die Zahlentheorie der Quaternionen (Berlin, Springer 1919), Lecture 6. There, however, what is essentially involved are principal ideals; the same holds for the works of Du Pasquier cited there in note 1.}

1. The \emph{right-sided module}, which, along with $M$, contains also $FM$ for an arbitrary polynomial $F$, and along with $M_1$ and $M_2$ contains also $M_1+M_2$.

2. The \emph{left-sided module}, which, along with $M$, contains also $MF$, and along with $M_1$ and $M_2$ contains also $M_1+M_2$.

We refer to these two kinds by the common name \emph{one-sided modules}; modules which are both right-sided and left-sided will be called \emph{two-sided}. In what follows we restrict ourselves to right-sided modules and note that the theory of the left-sided ones can be built up in exactly the same way.

We call two polynomials $F$ and $G$ congruent modulo $\mM$, written $F\equiv G(\mM)$, if their difference is divisible by $\mM$, that is, if it belongs to $\mM$. The totality of all polynomials congruent to a given polynomial forms a \emph{residue class} modulo $\mM$. (Residue classes are to be denoted by small German letters.) Since from
\[
 F_1\equiv F_2(\mM)\qquad\text{and}\qquad G_1\equiv G_2(\mM)
\]
it follows that $F_1+G_1\equiv F_2+G_2(\mM)$, the sum of two residue classes $\mx$ and $\my$ is defined and is again a residue class $\mx+\my$. By contrast, because the module is one-sided, one cannot speak of the product of two residue classes. Indeed, from
\[
 F_1=F_2+M_1,\qquad G_1=G_2+M_2
\]
there follows only
\[
 F_1G_1=F_2G_2+F_2M_2+M_1G_2+M_1M_2;
\]
but since $M_1G_2$ need not belong to $\mM$ when $M_1$ belongs to $\mM$, $F_1G_1$ and $F_2G_2$ need not be congruent in general. On the other hand, if $M_1=0$, hence $F_1=F_2=F$, one sees that $G_1\equiv G_2$ implies $FG_1\equiv FG_2$. Thus an arbitrary residue class $\my$ can be multiplied from the front by a polynomial $F$, and this gives a well-defined residue class $F\my=\mz$.

For computation with residue classes, the commutative and associative laws of addition hold, as does the law of unique reversal of addition; moreover, for multiplication of residue classes by polynomials and addition, the distributive law $F(\my_1+\my_2)=F\my_1+F\my_2$ holds, as follows directly from computing with polynomials modulo $\mM$.

The residue class to which the unit belongs will be called the unit class, or unit, and denoted $\me$. Then $F\me=\mx$ is the residue class to which the polynomial $F$ belongs. In contrast to the general residue classes, however, here the product $\mx\me$ is also well-defined and equal to $\mx$; for from
\[
 F_1=F_2+M_1,\qquad G_1=1+M_2
\]
and $M_1\cdot1=M_1$ it follows that
\[
 F_1G_1=F_2+M.
\]
The totality of the residue classes modulo $\mM$ is called the \emph{residue group of} $\mM$. A residue group therefore has the following properties: 1. Along with every residue class $\mx$ and an arbitrary polynomial $F$, it contains the residue class $F\mx$; and along with the residue classes $\mx_1$ and $\mx_2$ it contains the residue class $\mx_1+\mx_2$.\footnote{This property could, in a more general sense, be called the ``module property'' of the residue group.} 2. It possesses a unit $\me$, so that $F\me$ represents, for every polynomial $F$, the residue class $\mx$ of $F$. \emph{Thus, as} $F$ \emph{runs through all polynomials,} $F\me$ \emph{runs through the whole residue group.}

\subsection*{\S~4. Reducibility of the Residue Group into Two Subgroups and Its Relation to the Module.}

For one-sided modules, the concepts \emph{divisibility}, \emph{greatest common divisor}, and \emph{least common multiple} remain as in the two-sided case, as the following definitions show.

A module $\mM$ is said to be divisible by another module $\mN$ if from $M\equiv0(\mM)$ it follows also that $M\equiv0(\mN)$.

For two modules $\mM$ and $\mN$, the \emph{greatest common divisor} $(\mM,\mN)$ is defined as the totality of all polynomials representable in the form $M+N$, where $M\equiv0(\mM)$ and $N\equiv0(\mN)$; it is again a module. The same holds for several modules.

The totality of all polynomials divisible both by $\mM$ and by $\mN$ forms a module $[\mM,\mN]$, the \emph{least common multiple} of $\mM$ and $\mN$. This definition too applies correspondingly to several modules, indeed even to a set of modules of arbitrary cardinality.

Two modules are called relatively prime when their greatest common divisor is the unit module, that is, the totality of all polynomials.

We now start from a module $\mM$ which is representable as the least common multiple of two relatively prime modules $\mN$ and $\mL$, both assumed to be proper divisors:
\begin{equation}
 \mM=[\mN,\mL].
\tag{7}
\end{equation}
We investigate the residue group $\mG$ of $\mM$. By hypothesis,
\begin{equation}
 1\equiv N+L(\mM)
\tag{8}
\end{equation}
for two polynomials $N\equiv0(\mN)$ and $L\equiv0(\mL)$. From this we show
\begin{equation}
 \mL=(\mM,L),\qquad \mN=(\mM,N),
\tag{9}
\end{equation}
where, for example, $(\mM,L)$ means the module consisting of all polynomials of the form $M+FL$. In fact, first $(\mM,L)\equiv0(\mL)$. Conversely, from (8) it follows, for an arbitrary polynomial $F$, that
\begin{equation}
 F\equiv FN(\mL).
\tag{10}
\end{equation}
Now if $F\equiv0(\mL)$, then $FN\equiv0(\mL)$, hence $\equiv0(\mM)$, and consequently $F\equiv FL(\mM)$. In the same way one proves $\mN=(\mM,N)$. It follows incidentally, since $\mN$ and $\mL$ are proper divisors of $\mM$, that neither $L\equiv0$ nor $N\equiv0(\mM)$ can hold.

Let $\ma$ be the residue class of $N$ modulo $\mM$, and $\mb$ the residue class of $L$ modulo $\mM$. Then by (8)
\[
 \me=\ma+\mb\qquad(\ma\ne0,\ \mb\ne0).
\]
Since every residue class is of the form $F\me$, the distributive law gives, for every residue class,
\begin{equation}
 F\me=F\ma+F\mb\qquad(\ma\ne0,\ \mb\ne0).
\tag{11}
\end{equation}
This representation of an arbitrary residue class as a sum of two residue classes of the form $G\ma$ and $H\mb$ is, however, \emph{unique}. For from a relation $0=G\ma+H\mb$, if $K$ is a polynomial in the residue class $G\ma=-H\mb$, then plainly $K\equiv GN(\mM)$, $K\equiv-HL(\mM)$, hence $K\equiv0(\mN)$, $K\equiv0(\mL)$, and therefore $K\equiv0(\mM)$; that is, $G\ma=0$, $H\mb=0$.

Conversely, suppose now that relation (11) is fulfilled for every polynomial $F$, and indeed uniquely in the above sense. Let $N$ then be a polynomial from $\ma$ and $L$ a polynomial from $\mb$. We form the two modules (9) and assert (7) and (8). The latter follows from (11) for $F=1$. Furthermore, by the definition of $\mN$ and $\mL$, $\mM$ is a common multiple of these two modules, and indeed the least one. For from $F\equiv0(\mL)$ and $F\equiv0(\mN)$ there follows $F\equiv HL\equiv GN(\mM)$; hence $G\ma=H\mb$, $G\ma-H\mb=0$, and therefore, by uniqueness of the representation, $G\ma=0$, $H\mb=0$, $F\equiv0(\mM)$. It has therefore been shown that the \emph{representation of a module} $\mM$ \emph{as the least common multiple of two relatively prime modules is equivalent to the unique additive decomposition of the residue group} $\mG$ \emph{of} $\mM$, \emph{as given by formula} (11).

We now investigate the system $\mA$ of residue classes of the form $F\ma$, which we wish to relate to the residue group $\mbarA$ of $\mL$. By (11), every polynomial $F$ in $\mL$ belongs to the residue class $F\ma$ when $\bara$ is the residue class of $N$ modulo $\mL$, and hence the unit class of $\mbarA$. If we now associate the residue classes $F\ma$ in $\mA$ with the residue classes $F\bara$ in $\mbarA$, then to every residue class of $\mA$ there is assigned a residue class of $\mbarA$, in such a way that the sum of two residue classes of $\mA$ corresponds to the sum of the assigned residue classes of $\mbarA$, and the product of a residue class of $\mA$ with a polynomial corresponds to the product of the assigned residue class of $\mbarA$ with the same polynomial. This assignment is moreover one-to-one: from $F\ma=0$ follows $F\equiv0(\mL)$ by (11), hence $F\bara=0$; conversely, $F\bara=0$, hence $F\equiv0(\mL)$, also says by (10) that $FN\equiv0(\mL)$, hence $FN\equiv0(\mM)$, that is, $F\ma=0$.

This leads us to a fundamental concept, which we define as follows:

\emph{A one-to-one correspondence between two systems $\mA$ and $\mbarA$ of residue classes is called isomorphic if the sum of two residue classes of $\mA$ is assigned the sum of the corresponding residue classes of $\mbarA$, and the product of an arbitrary residue class of $\mA$ with a polynomial is assigned the product of the corresponding residue class of $\mbarA$ with the same polynomial.}

A subsystem $\mA$ of residue classes in $\mG$ which is isomorphic to the residue group of a module is called a \emph{subgroup} of $\mG$. As the proof above shows, the assignment here is specifically such that every residue class of $\mA$ arises from the corresponding residue class of $\mbarA$ by adjoining certain polynomials, namely all polynomials belonging to $\mL$ but not to $\mM$.

Similarly one shows that the totality of all residue classes of the form $F\mb$ forms a subgroup $\mB$ of $\mG$ which is isomorphic to the residue group $\mbarB$ of $\mN$.

A residue group $\mG$ whose residue classes admit a unique additive decomposition (11) will be called \emph{reducible} into the two subgroups $\mA$ and $\mB$, written $\mG=\mA+\mB$; a group which is not reducible is called \emph{irreducible}. We also occasionally call the associated modules reducible or irreducible. Thus the following holds.

\textbf{Theorem 1.}\footnote{Cf. Schmeidler, loc. cit., p. 39.} \emph{A module $\mM$ is representable as the least common multiple of two relatively prime proper divisors $\mL$ and $\mN$ if and only if its residue group $\mG$ is reducible into two subgroups $\mA$ and $\mB$. In this case $\mA$ and $\mB$ are respectively isomorphic to the residue groups of $\mL$ and $\mN$.}
\end{document}
